\documentclass[12pt]{article}
\usepackage[utf-8]{inputenc}
\usepackage{amsmath}
\usepackage{amssymb}

\title{Análisis de Propagación de Incertidumbre en el Logaritmo}
\author{}
\date{}

\begin{document}

\maketitle

\section*{Contexto}

El cálculo del flujo de calor del cuerpo está dado por:

\[
Q_{\text{cuerpo}} = \frac{\pi (T_{\text{cuerpo}} - T_{\text{ropa}}) \cdot k \cdot 2 \cdot h}{\ln(r_{\text{ext}}/r_{\text{int}})}
\]

donde el término del logaritmo aparece en el \textbf{denominador}.

\section*{Paso 1: Construcción del Ratio}

Primero se calcula el cociente entre los radios exterior e interior:

\begin{align*}
r_{\text{int}} &= 0.2 \pm 0.01 \text{ m} \\
r_{\text{ext}} &= 0.20350000000000001 \pm 0.01414213562373095 \text{ m} \\
\text{Ratio} &= \frac{r_{\text{ext}}}{r_{\text{int}}} = 1.0175 \pm 0.0500
\end{align*}

\section*{Paso 2: Propagación de Incertidumbre en el Logaritmo}

La propagación de incertidumbre en funciones se realiza mediante la derivada:

\[
\delta f(x) = \left| \frac{df}{dx} \right| \cdot \delta x
\]

Para el logaritmo natural:
\[
\frac{d(\ln x)}{dx} = \frac{1}{x}
\]

En nuestro caso:
\[
\frac{d(\ln)}{dx} = \frac{1}{1.0175} = 0.9828
\]

Por lo tanto, la incertidumbre en el logaritmo es:
\[
\delta(\ln(\text{ratio})) = 0.9828 \times 0.0500 = 0.0491
\]

\section*{Paso 3: Resultado del Logaritmo}

\[
\ln(1.0175) = 0.01735 \pm 0.0491
\]

\textbf{Incertidumbre relativa:} 
\[
\frac{\delta(\ln)}{\ln} \times 100\% = \frac{0.0491}{0.01735} \times 100\% = 283.29\%
\]

\section*{¿Por qué es tan grande la incertidumbre?}

Existen dos razones principales:

\subsection*{1. El logaritmo de un número cercano a 1 es muy pequeño}

Cuando el argumento del logaritmo está muy cercano a 1, el valor del logaritmo resultante es muy pequeño. En nuestro caso:
\[
\ln(1.0175) = 0.01735
\]

Una incertidumbre absoluta de $\delta(\ln) = 0.0491$ es enorme comparada con este valor pequeño, resultando en una incertidumbre relativa extremadamente grande.

\subsection*{2. La incertidumbre del espesor de la ropa es desproporcionada}

El espesor de la ropa es medido como:
\[
\text{espesor}_{\text{ropa}} = 0.0035 \pm 0.01 \text{ m}
\]

La incertidumbre relativa es:
\[
\frac{\delta(\text{espesor})}{\text{espesor}} = \frac{0.01}{0.0035} = 285.7\%
\]

Esta incertidumbre relativa extremadamente grande se propaga directamente al radio exterior y, posteriormente, al ratio $r_{\text{ext}}/r_{\text{int}}$, amplificando la incertidumbre en el logaritmo.

\section*{Impacto en el Flujo del Cuerpo}

Dado que el logaritmo aparece en el denominador de la fórmula para $Q_{\text{cuerpo}}$, la gran incertidumbre relativa del logaritmo genera una incertidumbre aún más pronunciada en el resultado final:

\[
Q_{\text{cuerpo}} = 3549.65 \pm 10056.28 \text{ kcal/día}
\]

Con una incertidumbre relativa de aproximadamente $284\%$, lo que indica que esta magnitud tiene una precisión muy baja debido principalmente a la propagación de la incertidumbre a través del logaritmo.

\end{document}
